\section{Secure Mobile Pairing and Mother-Brain Continuity}

The unified Pilot application treats a phone as a private cryptographic key and user interface,
not as an independent source of authority.  The user-selected mother brain remains authoritative
for identity, memory, capabilities, approvals, execution and receipts.  This stage reuses Pilot's
existing private-identity, local TLS, recovery and shared-screen systems and adds one coherent
phone-to-mother trust protocol around them.

\subsection{Security objectives}

The pairing design has five primary objectives:
\begin{enumerate}
  \item connect a phone without exposing a mother-brain password or private key;
  \item allow the phone to prove possession of its own key;
  \item allow the phone to verify that it is connecting to the intended mother brain;
  \item issue only explicit, time-bounded capabilities; and
  \item make loss, theft, expiry and revocation fail closed across messaging, approvals and shared
  television sessions.
\end{enumerate}

The protocol does not treat knowledge of a URL, QR code or six-digit number as durable authority.
Durable trust is established only after a signed challenge and is represented by a mother-signed
phone certificate plus a short-lived capability lease.

\subsection{Signed mother discovery}

The local HTTPS service publishes a short-lived descriptor at
\texttt{/.well-known/pilot-mother}.  The descriptor contains the mother identifier, Ed25519 public
key, key fingerprint, HTTPS endpoint, household certificate-authority fingerprint and bounded
timestamps.  The whole descriptor is signed by the mother.  The phone verifies the signature,
recomputes the fingerprint, rejects expired descriptors and accepts only a path-free HTTPS
endpoint.  Altering either the endpoint or key therefore invalidates verification.

Pilot also derives a small set of human-readable trust words from the mother public key.  The
same words are displayed on the phone and trusted household surface so that a user can detect a
substituted QR code or unexpected mother brain without reading hexadecimal fingerprints.

\subsection{Pairing ceremony}

The connection sequence is:
\begin{enumerate}
  \item The phone discovers or scans the signed mother descriptor and verifies it locally.
  \item The phone creates an Ed25519 device key.  In the web demonstrator the temporary exported
  key bytes are erased and the private key is re-imported as non-extractable before storage in
  IndexedDB.  Native applications will use the platform key store where available.
  \item The phone requests a five-minute, single-use challenge containing its public key, selected
  persona and requested capability scopes.
  \item The mother signs the challenge.  A random six-digit confirmation code is shown through a
  trusted local presenter on the mother side.  The HTTPS response deliberately never contains the
  code.
  \item The user enters the code on the phone, which signs the canonical challenge payload.
  \item The mother verifies the code, expiry, attempt count and phone signature, then records the
  phone in the existing Production Private Identity registry.
  \item The mother returns a signed phone certificate and an eight-hour capability lease.
\end{enumerate}

A challenge can be consumed only once.  Five incorrect confirmation attempts lock it.  A newer
challenge for the same persona and phone supersedes an older pending challenge.  This bounds both
brute-force attempts and accidental duplicate enrollment.

\subsection{Certificates, leases and reconnection}

The phone certificate binds the device identifier, persona, public key, mother, household TLS
fingerprint and validity interval.  The separate capability lease binds that certificate to an
ordered list of allowed scopes and expires after eight hours by default, with a hard maximum of
twenty-four hours.

Lease renewal requires a fresh nonce and a phone signature.  The previous lease is rotated rather
than silently reused, replayed renewal nonces are rejected, and a renewal may retain or reduce
scope but may not add a new permission.  Every protected request validates both mother signatures,
both expiry times, the stored active lease and the authoritative private-identity device status.

Direct trusted-LAN access is preferred at home.  The interface uses ordinary descriptions such as
``Connected securely'', ``Needs reconnecting'', ``Offline'' and ``Phone removed'' instead of
exposing certificate terminology to ordinary users.

\subsection{Revocation and recovery}

There is one authoritative device lifecycle.  Revoking a lost phone through Pilot's private
identity settings marks the device revoked, invalidates every active mobile lease and terminates
its shared-screen authority.  The pairing layer checks that authoritative status on validation,
renewal and connection-state queries, so a stale local certificate cannot restore access.

Existing one-time recovery codes and explicit identity recovery remain the recovery mechanism.
Recovery does not copy an old private key onto a replacement device.  The replacement completes a
new possession ceremony, and the lost device remains revoked.

\subsection{Network and implementation controls}

The local pairing API requires TLS~1.2 or newer, exact-origin allowlisting, JSON requests below
32~KiB, no-store responses, strict transport security, frame denial and a bounded per-address
request rate.  Private mother and certificate-authority keys never enter the HTTP response.
Persisted pairing state and identity state use owner-only file permissions.

Automated adversarial tests cover endpoint and public-key substitution, incorrect-code locking,
single-use challenges, phone-possession signatures, certificate and lease mismatch, scope
expansion, renewal replay, direct identity-registry revocation, public inventory redaction and a
real local HTTPS round trip whose phone response is inspected to prove that it does not contain the
confirmation code.

\subsection{Completion boundary and remaining continuity work}

This stage completes local mother discovery, the pairing ceremony, certificate and lease lifecycle,
non-technical connection states, device inventory, recovery integration and lost-phone revocation.
It does not claim that remote continuity is complete.  The following items remain:
\begin{itemize}
  \item an encrypted peer-to-peer route outside the home;
  \item an end-to-end encrypted relay fallback whose operator cannot read application content;
  \item automatic Wi-Fi to mobile-data roaming with bounded retry and route health;
  \item resumable subscriptions with message and action duplicate suppression;
  \item hardware-backed native iOS and Android keys and biometric approval; and
  \item field exercises across home Wi-Fi, external Wi-Fi, mobile data, relay failure and mother
  migration.
\end{itemize}

These remaining network transports must reuse the certificate, lease and revocation rules above.
They must not introduce a second identity authority or turn a relay account into authority over a
user's Pilot.
